\documentclass[paper=a4, fontsize=11pt, BCOR=13mm, DIV=13, headinclude, toc=index, toc=bibliography, english, twoside]{scrreprt}
% Die verwendete Dokumentenklasse ist scrreprt. Die verwendeten Optionen sind:
%
% paper=a4              Papier ist a4.
% fontsize=11pt         Schrifgröße ist 11.
% DIV=13                Das Papier wird in d viele Spalten und d' viele Zeilen eingeteilt. Die Werte werden aus DIV berechnet.
% BCOR=1cm              Definiert den Patz, der auf der Innenseite beim Binden verloren geht.
% headinclude           sorgt dafür, dass genug Platz für die Header vorhanden ist.
% toc=index             legt im Inhaltsverzeichnis einen Eintrag für das Stichwortverzeichnis an.
% toc=bibliography      legt im Inhaltsverzeichnis einen Eintrag für das Literaturverzeichnis an.
% english               englische Worte wie "Chapter" und "References".
% twoside               Beidseitiges Dokument, wie in einem Buch.


%If you don't want this fancyheaders comment out lines 17 to 24
% Fancyheader
\usepackage{fancyhdr}                   % Wie der Name schon sagt, um fancy Header zu generieren.
\pagestyle{fancy}                       % Fancy Header sollen angezeigt werden
\renewcommand{\sectionmark}[1]{\markright{\thesection.\ #1}{}}    % Verhindert dass rightmark ausschließlich Grußbuchstaben benutzt
\fancyhead[LE,RO]{\rightmark}           % Links bei geraden und rechts bei ungeraden Seitenzahlen soll der Name der Section stehen.
\fancyhead[LO,RE]{}                     % Links bei ungeraden und rechts bei geraden Seitenzahlen soll nichts stehen.
\fancyfoot[C]{}                         % Keine mittigen Seitenzahlen
\fancyfoot[LE,RO]{\thepage}             % Seitenzahlen unten in die jeweilige äußere Ecke





\setcounter{secnumdepth}{3}     % Nummerierungstiefe (chapter, section, subsection, ...).
\setcounter{tocdepth}{3}        % Nummerierungstiefe im Inhaltsverzeichn is.

\usepackage[linesnumbered,ruled,vlined]{algorithm2e}    % Algorithmen setzen.
\usepackage{amsmath,amssymb,amsthm,amsfonts,amsbsy,latexsym}    % "Notwendige" AMS-Math Pakete.
\usepackage{array}                      % Bessere Tabellen.
\renewcommand{\arraystretch}{1.15}      % Tabellen bekommen ein wenig mehr Platz.
\usepackage{bbm}                        % Dicke 1.
\usepackage[backend=biber, style=alphabetic]{biblatex}  % Gute Erweiterung zu bibtex, Wird für Referenzen benutzt.
\bibliography{thesis}   % Die verwendeten Referenzen (.bib-Datei)
\usepackage[hypcap]{caption}            % Damit Hyperrefs bei der figure-Umgebung auf die Figure zeigt statt auf die Caption.
\usepackage[nodayofweek]{datetime}                   % Um \today einzustellen.
\newdateformat{mydate}{\THEDAY{}th \monthname{} \THEYEAR{}}
\usepackage{diagbox}                    % Diagonale in Tabellen.
\usepackage{enumitem}                   % Zum Ändern der Nummerierungsumgenung 'enumerate'
\setlist[enumerate,1]{label=(\roman*)}  % Aufzählungen sind vom Typ 'Klammer auf; kleine römische Zahl; Klammer zu'
\usepackage[T1]{fontenc}                % Bessere Schrift
\usepackage{ifthen}                     % Zum checken ob Parameter leer sind.
\usepackage{lmodern}                    % Bessere Schrift
\usepackage{listings}                   % Code Listings.
\usepackage{mathtools}                  % Subscript unter Summen behandeln. Der Befehl lautet \mathclap.
\usepackage{makeidx}                    % Stichwortverzeichnis.
\makeindex                              % Stichwortverzeichnis erstellen.
\renewcommand{\indexname}{Index}        % Name des Index definieren.
\usepackage{multirow}                   % In Tabellen mehrere Zeilen zu einer machen.
\usepackage{rotating}                   % Um Figures zu drehen.
\usepackage{scrhack}                    % Verbessert die Zusammenarbeit von KOMA mit anderen Paketen (z.B, listing).
\usepackage{stackrel}                   % Symbole übereinander stapeln.
\usepackage[dvipsnames]{xcolor}         % Gefärbter Text und so.
\usepackage{tikz}                       % Graphen und kommutative Diagramme. Muss nach xcolor eingebunden werden.
\usepackage{tikz-cd}                    % Kommutative Diagramme.
\usepackage{comment}                    % Multiline comments

\graphicspath{{pictures/}}              % Pfad in dem die mit Inkscape erstellen Bilder liegen (relativ zum Hauptverzeichnis).

% Workaround, damit keine unnötigen Leerzeichen entstehen.
\let\oldindexdefn\index
\renewcommand*{\index}[1]{\oldindexdefn{#1}\ignorespaces}
\let\oldlabeldefn\label
\renewcommand*{\label}[1]{\oldlabeldefn{#1}\ignorespaces}

% Workaround, Linebreak nach ldots erlaubt.
\newcommand{\origldots}{}
\let\origldots\ldots
\renewcommand{\ldots}{\allowbreak\origldots}

% to get upright greek letters
\usepackage{upgreek}

% To break urls correctly in bibliography
\setcounter{biburllcpenalty}{7000}
\setcounter{biburlucpenalty}{8000}

% to get rid of an error
\usepackage{bookmark}

%colors for lean syntax
\definecolor{keywordcolor}{rgb}{0.7, 0.1, 0.1}   % red
\definecolor{tacticcolor}{rgb}{0.0, 0.1, 0.6}    % blue
\definecolor{commentcolor}{rgb}{0.4, 0.4, 0.4}   % grey
\definecolor{symbolcolor}{rgb}{0.0, 0.1, 0.6}    % blue
\definecolor{sortcolor}{rgb}{0.1, 0.5, 0.1}      % green
\definecolor{attributecolor}{rgb}{0.7, 0.1, 0.1} % red

% To display lean code
\def\lstlanguagefiles{lstlean.tex}
% set default language
\lstset{language=lean, extendedchars=true}

%to get code in footnotes write \cprotect\footnote 
\usepackage{cprotect}


% Symbolverzeichnis
\usepackage[intoc, english]{nomencl}     % Symbolverzeichnis.
% intoc                 die Symbolliste in das Inhaltsverzeichnis aufnehmen.
% english               englische Worte wie "Seite".
\renewcommand{\nomname}{Symbol Index}   % Definiert die Überschrift des Symbolverzeichnises.
\renewcommand{\nomlabelwidth}{80pt}     % Platz der einem Symbol gegönnt wird.
\newcommand{\symbolindex}[4][]{{\nomenclature[#1]{#2}{#3\ifthenelse{\equal{#4}{}}{}{ -- #4}\nomnorefpage}}\ignorespaces}    % Verbesserte Version von "\nomenclature". Erzeugt Symbol Beschreibung - Referenz.
\renewcommand*{\nompreamble}{\markright{\nomname}}    % Workaround: Fancyhdr schreibt im Symbolverzeichnis sonst den Namen des leztztes Kapitels.
\makenomenclature                       % Symbolverzeichnis erstellen.

\usepackage{MA_Titlepage}

% to get the external link symbol
\usepackage{fontawesome}

% Anklickbare Referenzen (letztes eingebundenes Paket)
\usepackage{hyperref}                   % Referenzen innerhalb des Dokuments anklickbar machen. Achtung: Muss das letzte Paket im Präambel sein.
\hypersetup{                            % Optionen von hyperref Einstellen.
    colorlinks=true,                    % gefärbte Links an Stelle von Boxen.
    linkcolor=black,                     % Farbe interner Links.
    citecolor=black,                     % Farbe von Referenzen.
    urlcolor=black                       % Farbe von Internetlinks.
}



%Name of the author of the thesis 
\authornew{Hannah Scholz}
%Date of birth of the Author
\geburtsdatum{\formatdate{21}{7}{2003}}
%Place of Birth
\geburtsort{Bonn, Germany}
%Date of submission of the thesis
\date{\formatdate{22}{8}{2024}}

%Name of the Advisor
% z.B.: Prof. Dr. Peter Koepke
\betreuer{Advisor: Dr. Michael Rothgang}
%name of the second advisor of the thesis
\zweitgutachter{Second Advisor: Prof. Dr. Floris van Doorn}

%Name of the Insitute of the advisor
%z.B.: Mathematisches Institut
%\institut{Institut XYZ}
\institut{Mathematisches Institut}
%\institut{Institut f\"ur Angewandte Mathematik}
%\institut{Institut f\"ur Numerische Simulation}
%\institut{Forschungsinstitut f\"ur Diskrete Mathematik}
%Title of the thesis 
\title{????}
%Do not change!
\ausarbeitungstyp{Master's Thesis  Mathematics}

%       Theoreme
\theoremstyle{definition}               % Name: dick            Text: normal.
\newtheorem{defi}{Definition}[section]  % Der Zähler ist defi = Sectionzähler.1 . Sectionzähler soll bei Benutzung von defi nicht erhöht werden.
\newtheorem*{defi*}{Definition}
\newtheorem{example}[defi]{Example}
\newtheorem{notation}[defi]{Notation}
\newtheorem{rem}[defi]{Remark}
\newtheorem{defcor}[defi]{Definition/Corollary}
\newtheorem{defprop}[defi]{Definition/Proposition}
\newtheorem{defthm}[defi]{Definition/Theorem}

\theoremstyle{plain}
\newtheorem*{conj}{Conjecture}
\newtheorem{cor}[defi]{Corollary}
\newtheorem{lem}[defi]{Lemma}
\newtheorem{prop}[defi]{Proposition}
\newtheorem*{prop*}{Proposition}
\newtheorem{thm}[defi]{Theorem}
\newtheorem*{thm*}{Theorem}


% Claim environment (nestable once)
\begin{comment}
    The syntax is: 
    \begin{claim}
        Claim statement
        \begin{proof}
            Proof of claim.
        \end{proof}
    \end{claim}
\end{comment}

\newlist{Claim}{description}{2}
\setlist[Claim]{labelindent=2em,leftmargin=*}
\newif\ifInsideClaim
\newcounter{claim}[defi]
\newcounter{cclaim}[claim]
\renewcommand\theclaim{\arabic{claim}}
\renewcommand\thecclaim{\arabic{claim}.\arabic{cclaim}}
\let\originalqedsymbol\qedsymbol
\newenvironment{claim}{%
  % disable qed symbol if there is no star
  \let\qedsymbol\relax%
  \ifInsideClaim% we have a nested environment
    \refstepcounter{cclaim}%
    \let\theclaimcounter\thecclaim%
  \else%
    \refstepcounter{claim}%
    \let\theclaimcounter\theclaim%
    \InsideClaimtrue%
  \fi%
  \Claim\item[\textbf{Claim \theclaimcounter:}]%
}{\endClaim\InsideClaimfalse\let\qedsymbol\originalqedsymbol}

%       Makros
\input{thesis_macros}


\begin{document}
% Titelseite
\maketitle              % Titelseite ausgeben
\setcounter{page}{3}    % Die Titelseite und die darauffolgende leere Seite sollen gefälligst Seite 1 und 2 sein.
\tableofcontents        % Inhaltsverzeichnis ausgeben

%mathematics of CW-complexes
\cleardoublepage 
\chapter{The mathematics of collar neighborhoods}

\section{Prerequisites}
\label{sec:prereq}

\begin{lem}[Smooth Urysohn's lemma, in mathlib I think]
    Let $M$ be a $\fC^k$-manifold and $A$ and $B$ two disjoint closed sets in $M$. 
    Then there is a $C^k$-function $f \colon M \to [0, 1]$ such that $f$ is $0$ on $A$ and  $1$ on $B$.
\end{lem}

\begin{rem}
    Again, this requires a definition of manifold that makes it metrizable. 
\end{rem}
\section{Basics about collar neighbourhoods}
\label{sec:mathsdef}

\begin{defi}
    A \emph{topological manifold} is a second-countable Hausdorff space that is locally Euclidean.
\end{defi}

\begin{rem}
    Apparently, there are some subtleties with collars of non-metrizable manifolds but I assume we don't care. (At least with the above definition our manifolds are always metrizable.)
\end{rem}

\begin{defi}
    Let $M$ and $N$ be two $C^k$-manifolds. 
    Let $f \colon M \to N$. 
    \begin{enumerate}
        \item $f$ is a $C^k$-diffeomorphism if it is bijective and both $f$ and its inverse functions are $k$-times continuously differentiable. 
        \item $f$ is a $C^k$-embedding if it is a $C^k$-diffeomorphism onto its image.
    \end{enumerate}
\end{defi}

\begin{defi}[Due to \cite{Connelly1971}]\label{defi:collared}
    Let $M$ be a $\fC^k$-manifold (with boundary) and let $B \subseteq M$ be closed. 
    \begin{enumerate}\label{defi:collared:local}
        \item $B$ is \emph{locally collared} if it is covered by sets $U \subseteq B$ which are open in $B$ such that there are $\fC^k$-embeddings $\varphi_U \colon \closure{U} \times [0, 1] \to M$ with 
        \begin{enumerate}
            \item $\preim{\varphi_U}{B} = \closure{U} \times \{ 0\}$.
            \item For all $x \in \closure{U}$, $\varphi_U (x, 0) = x$.
            \item $\varphi_U (\closure{U} \times [0, 1])$ is closed. 
            \item $\varphi_U (U \times [0, 1)$ is open.  
        \end{enumerate}
        \item $B$ is \emph{collared} if $B$ already fulfils the conditions given above. In this case we call the image of $B \times [0, 1)$ under the corresponding map $\varphi_B \colon B \times [0, 1] \to M$ a \emph{collar neighbourhood} of $B$.
    \end{enumerate}
\end{defi}

\begin{rem}
    This definition is a lot more restrictive than other definition which don't consider the closure or the closed interval. I think if this one doesn't work, I can just take the other one. 
\end{rem}

\begin{rem}
    For $k=0$ we can even extend this definition to topological spaces instead of just defining it for topological manifolds. 
\end{rem}

\begin{thm}[Due to \cite{Connelly1971}]
    Let $X$ be a Hausdorff space and $B \subseteq X$ be a closed subset. 
    If $B$ is locally collared in $X$, then it is collared.
\end{thm}

\begin{proof}
    Let $U_1, \dots , U_n$ be a cover of $B$ that witnesses that $B$ is locally collared. 
    Let $\varphi_i \colon \closure{U_i} \times [0, 1] \to X$ be the corresponding embeddings. 
    $B$ is normal because it is compact and Hausdorff. 
    By the shrinking lemma (already in mathlib) there is a new open cover $V_1, \dots, V_n$ of $B$ such that $\closure{V_i} \subseteq U_i$ for all $1 \leq i \leq n$. 
    Define 
    \begin{equation*}
        X^+ \coloneq (X \amalg B \times [-1, 0]) / x \sim (x, 0).
    \end{equation*}
    Clearly, $X^+$ is collared. So it suffices to show that there is a homeomorphism $G \colon X \to X^+$ s.t.\ $G(x) = (x, -1)$ for all $x \in B$. 
    Let $\Phi_i \colon \closure{U_i} \times [-1, 1] \to X^+$ be defined by 
    \begin{equation*}
        \Phi_i (x) = \begin{cases*}
            \phi_i(x) & \text{if } $x \in \closure{U_i} \times [0, 1]$ \\
            x & \text{if } $x \in \closure{U_i} \times [-1, 0]$. \\
        \end{cases*}
    \end{equation*}
    We now inductively define maps $f_i \colon B \to [0, 1]$ and embeddings $g_i \colon X \to X^+$ for all $0 \leq i \leq n$ such that 
    \begin{enumerate}
        \item $f_i(x) = -1$ if $x \in \bigcup_{j \leq i} \closure{V_j}$
        \item $g_i(x) = (x, f_i(x))$ if $x \in B$
        \item $g_i(X) = X \cup \{ (x, t) \mid t \geq f_i(x)\}$
    \end{enumerate}
    Then $g_n$ will be the desired homeomorphism $G$ from above.
    Set $f_0$ to be constant $0$. Then we can choose $g_0$ to be the inclusion and trivially $g_0(X) = X = X \cup X \times \{0\} = X \cup \{ (x, t) \mid t \geq 0\}$. 
    
    Assume that $g_{i-1}$ has been defined. Using Urysohn's Lemma, choose a continuous function $\lambda_i \colon \closure{U_i} \to [0, 1]$ such that it is $0$ on $\closure{U_i} \setminus U_i$ and $1$ on $\closure{V_i}$. 
    Given $x \in B$ set $a(x) \coloneq f_{i-1}(x)$ and $b(x) \coloneq (1 - \lambda_i(x))f_{i-1}(x) + \lambda_{i}(x) (-1)$.
    Now let 
    \begin{equation*}
        s(x, \_) \colon [a(x), 1] \to [b(x), 1],\quad t \mapsto \frac{b(x)-1}{a(x)-1}(t-1) + 1.
    \end{equation*}
    Now set
    \begin{equation*}
        \psi_i \colon \preim{\Phi_i}{g_{i-1}(X)} \to \closure{U_i} \times [-1,1], \quad (x, t) \mapsto (x, s(x, t)).
    \end{equation*}
    This allows us to define 
    \begin{equation*}
        \Psi_i \colon g_{i-1}(X) \to X^+, \quad x \mapsto \begin{cases*}
            \Phi_i (\psi_i (\inv{\Phi_i}(x))) & if  $x \in g_i(X) \cap \Phi_i(\closure{U_i} \times [-1, 1])$ \\ 
            x & \text{otherwise} \\
        \end{cases*}
    \end{equation*}
    \begin{claim}
        $\Psi_i$ is continuous. 
    \end{claim}
    \begin{proof}
        It suffices to show that $\psi_i$ is continuous. 
        We need to show that $s \colon B \times [a(x), 1] \to [b(x), 1]$ is continuous. 

        Ahhh this is ugly. I somehow would prefer to just get going with this in Lean.
    \end{proof}
\end{proof}

I want to try to separate out some parts here: 

\begin{defi}
    Let $X$ be a $\fC^k$-manifold and $B \subseteq X$ a closed subset. 
    We define the topological space
    \begin{equation*}
        X^+_B \coloneq (X \amalg B \times [-1, 0]) / x \sim (x, 0).
    \end{equation*}
\end{defi}

\begin{lem}
    Let $X$ be an $n$-dimensional $\fC^k$-manifold and $B \subseteq X$ a closed subset. 
    Then $X^+_B$ is an $n$-dimensional $\fC^k$-manifold. 
\end{lem}
\begin{proof}
    We first need to show that $X^+_B$ is a Hausdorff space. I am not sure if this is best done manually or if there is some cool lemma like: a pushout along closed embeddings preserves Hausdorff spaces. (Is this true?)

    Then we need to show that $X^+_B$ is second countable. Again, there should be some general lemma about closed embeddings and pushouts. 

    Now we need to pick an atlas for $X^+_B$.
    Let $\fA_X$ be a $\fC^k$-atlas for $X$, $\fA_B$ be a $\fC^k$-atlas for $B$ and $\fA_{[-1, 0]}$ be a $\fC^k$-atlas for $[-1, 0]$. 
    We define 
    \begin{align*}
        \fA \coloneq &\{(U, \varphi) \in \fA_X \mid U \cap B = \varnothing\} \\
        &\cup \{(U \times I, \varphi \times \psi) \mid (U, \varphi) \in \fA_B, (I, \psi) \in \fA_{[-1, 0]}, 0 \notin I\} \\
        &\cup \{(U \cup ((U \cap B) \times I), \varphi^+_I) \mid (U, \varphi) \in \fA, (I, \psi) \in \fA_{[-1, 0]}, 0 \in I\}
    \end{align*}
    where we define $\varphi^+_I$ as follows. 
    Let $(U, \varphi)$ and $(I, \psi)$ as above. 
    Set 
    \begin{equation*}
        \varphi^+_I \colon U \cup ((U \cap B) \times I) \to \bR^n, x \mapsto \begin{cases*}
            \varphi (x) & $x \in U$ \\ 
            (\varphi(y), z) & $x = (y, z) \in (U \cap B) \times I$ \\
        \end{cases*}
    \end{equation*}
    For $x \in B$ we have $\varphi(x)^+_I = \varphi(x) = (\varphi(x), 0) = \varphi^+_I(x, 0)$, so $\varphi^+_I$ is continuous. 

    \textbf{It seems to me that this might not be $\fC^k$ at all. I don't know of you can even glue this without basically already having the vector field available and then one might as well do the other proof. Maybe a local vector field is enough? Not sure :(}

    Now you would need to verify that this is a $\fC^k$ atlas. 
    Possibly using a version of the manifold chart lemma from Lee. 
    I will only try to show that the transition maps are $\fC^k$. 
    The rest should be easier. 
    The only interesting intersections are those of two charts of the form $(U \cup ((U \cap B) \times I), \varphi^+_I)$ and $(V, \cup ((V \cap B) \times J), \psi^+_J)$. 
    We need to check that the transition map is $\fC^k$ on $U \cap V \cap B$. 
    This is clear for all combinations of partial derivatives that do not include the $n$-th one. 
\end{proof}

\begin{lem}
    Let $X$ be an $n$-dimensional $\fC^k$-manifold and $B \subseteq X$ a closed subset. 
    Then $X^+_B$ is the pushout in the category of $\fC^k$-manifolds. 
\end{lem}

\begin{rem}
    It seems that this is not at all trivial and in fact much of the issue with gluing manually along boundaries.
    I am not sure if we can even guarantee the above lemma.
    The above lemma would hopefully mean that one never has to look at the explicit construction again.
\end{rem}


\begin{defi}
    Let $X$ be a Hausdorff space and $B \subseteq X$ a closed subset.
    Let $U \subseteq B$ be an open subset together with a 
\end{defi}

We would like the following theorem: 

\begin{thm}
    Let $M$ be a $\fC^k$-manifold (with boundary) and $B \subseteq X$ a closed subset. 
    If $B$ is locally collared in $M$, then it is collared.
\end{thm}

Then we will need the following lemma: 

\begin{lem}
    Let $M$ be an $n$-dimensional $\fC^k$-manifold with boundary. 
    Then $\boundary M$ is locally collared.
\end{lem}

\begin{proof}
    Let $p \in \boundary M$. Then $p$ is contained in a regular half-coordinate ball $U$, i.e.\ there is a coordinate half-ball $U' \supseteq \closure{U}$ and a coordinate map $\psi \colon U' \to \bH^n$ such that for some $r' > r > 0$ we have $\psi(U) = B_r(0) \cap \bH^n$, $\psi (\closure{U}) = \closure{B_r(0)} \cap \bH^n$ and $\psi (U') = B_{r'}(0) \cap \bH^n$. 

    Now pick an open neighbourhood $V$ of $0 \in \boundary \bH^n$ and $l > 0$ s.t. $V \times [0, l] \subseteq B_r(0)$.

    Considering $\closure{V}$ and $V$ as submanifolds of $\bH^n$ and $\preim{\psi}{\closure{V}} = \closure{\preim{\psi}{V}}$ and $\preim{\psi}{V}$ as submanifolds. 
    Then 
    \begin{equation*}
        \restrict{\psi}{\closure{\preim{\psi}{V}}} \colon \closure{\preim{\psi}{V}} \to \closure{V}
    \end{equation*}
    and 
    \begin{equation*}
        \restrict{\psi}{\preim{\psi}{V}} \colon \preim{\psi}{V} \to V
    \end{equation*}
    are $\fC^k$-diffeomorphisms. 
    So it suffices to show that $\restrict{\inv{\psi}}{\closure{V} \times [0, l]} \colon \closure{V} \times [0, l] \to M$ fulfils the conditions given in Definition \ref{defi:collared}. 

    I think this should work... 
    I think it would be nicer to do this directly in Lean and then later write up the maths proof. 
\end{proof}

which will then allow us to derive

\begin{thm}[Collar neighbourhood theorem]
    Let $M$ be a $C^k$-manifold with boundary, then $\boundary M$ is collared. 
\end{thm}

%\begin{defi}
%    Let $M$ be a $\fC^k$-manifold with boundary and let $U$ be a neighbourhood of the boundary $\boundary M$. 
%    $U$ is called a \emph{collar neighbourhood} if there is a $\fC^k$-diffeomorphism $f \colon \boundary M \times [0, 1) \to U$ s.t.\ for all $x \in \boundary M$, $f(x, 0) = x$. 
%\end{defi}


% Symbolverzeichnis
\cleardoublepage        % Auch diese sollen auf der rechten Seite beginnen
\printnomenclature      % Symbolverzeichnis ausgeben

% Referenzen
\nocite{*}              % Alle Einträge der Bib-Datei sollen in die Referenzen
\cleardoublepage        % Auch diese sollen auf der rechten Seite beginnen
\printbibliography      % Bibliographie ausgeben.

\end{document}